\section{Processors and Controllers}
\subsection{Microprocessor}
The microprocessor is \it{just} the computing unit and has to be complemented by different types of memory, I/O interfaces and other system components. They are usually fast but therefore expensive.

\subsection{Microcontroller}
A microcontroller (\f{\mu C}) is a System-on-Chip (SoC), which contains all basic components such as memory, timers, and various I/O interfaces (I\f{^2}C, SPI, UART). In comparison to microprocessors they are slower, but also cheaper and more energy efficient.

Classically they are used, when control operations are more frequent than data intensive computing (e.g. in multimedia).
\newline

To use a \f{\mu C}, the ES needs to fulfill multiple efficiency requirements: Energy efficiency (especially when using batteries), memory efficiency (thus code and architecture efficiency), run-time efficiency (processing efficiency for real-time use critical)